\chapter{An Extended Example}
\label{chapter-examples}

	The strong head normalizing semantics of THOR are unique enough to
warrant providing an example of how it can be used effectively. This
chapter will demonstrate how THOR can be used to constructively program
numerical functions.  The sine example given in this chapter is adapted
from a paper by Roylance \cite{roylance}, as 
is much of the accompanying discussion. 
The original work was presented in Scheme, but THOR provides a more
natural vehicle for this style of programming.

\section{Constructive Programming}

	Constructive programming is an `indirect' programming paradigm in
which you don't merely write a procedure for computing values, but you
write a procedure for constructing a procedure that computes values.
Constructive programming is particularly useful for numerical programs,
because they can often be written to closely match their mathematical
basis.

	Constructive programs are generally longer and often more difficult
to write than conventional programs.  Then why bother?  Because they
offer some significant advantages.  The
clarity and believability of a program is enhanced because much of its
theoretical foundation is included explicitly in the code.  This also
helps to
reduce the chance of programming errors.  In addition, a single
constructive program specifies a whole family of procedures whose
characteristics, such as accuracy or dimension, can be changed by trivial
modifications to input parameters. 


	In many languages, simply executing constructive programs as
they are written (without any kind of preprocessing) would be grossly
inefficient.  In THOR it is possible (using the {\tt RDEFINE} special
form) to `reduce away' much of the constructive description, generating
more efficient code.  
	\footnote{The term {\em abstract interpretation} is often used to
	describe the symbolic, meta-execution of a program without data.  However,
	this term is not appropriate to THOR, because there is nothing
	`meta' about what is going on here.  The same reduction mechanism is
	used to reduce programs without data as is used for programs with
	data.}

	As an example of using {\tt RDEFINE} and strong head normalization in
THOR, the rest of this chapter is
devoted to constructively programming the trigonometric sine function as
a simple Taylor series expansion.  In the process of constructing an
expression for the
sine function, we will develop several general purpose higher order
expressions that are useful in other constructive programming projects.

	I do not claim simple Taylor series 
expansions are the best way to compute the sine function, but most 
people are familiar
with them.  The same techniques could be used with more sophisticated 
methods of
computing sine, such as Chebyshev economization.


\section{Approximating the Sine Function}

	The trigonometric function $\sin(x)$ can be represented by the
Taylor Series
\[\sum_{i=0}^{\infty} (-1)^{i} \frac{1}{(2i+1)!} x^{2i+1} \]
This particular series is absolutely convergent for all values of $x$. 
We shall use a finite approximation of this infinite series as the basis of our 
sine function.  

	Before constructing an approximation to this series,
let us first take a moment to think about the nature of numerical
approximations in general, and then about this one in particular. 
Approximations are most effective if they only cover small
intervals, so we would like to find ways to reduce the interval size.  This
is easy to do with the sine function because it is periodic, with period
$2\pi$. If the periodic interval is chosen to be [$-\pi/2, 3\pi/2$], 
the interval size can be reduced even further because this particular 
interval is symmetrical about $\pi/2$. 

	Periodicity can be exploited to reduce interval size in many numerical
programs, so it is worthwhile to encapsulate this notion into a higher
order expression:
\begin{tt}\begin{tabbing}
(de\=fine period-maker\\
\>	(la\=mbda (fcn a period x)\\
\>\>	(le\=t ((cycles (floor (/ (- x a) period))))\\
\>\>\>	(fcn (- x (* cycles period))))))
\end{tabbing}\end{tt}
This expression takes a function, {\tt fcn}, of the variable {\tt x},
that is defined over the interval [{\tt a}, {\tt a+period}] and reduces
to an expression that replicates {\tt fcn} over the rest of the real
number line.


	Symmetry can also be exploited to reduce interval size in many
programs.  Reflection of a function about an axis of symmetry is
encapsulated by this higher order expression:
\begin{tt}\begin{tabbing}
(de\=fine reflection-maker\\
\>	(la\=mbda (fcn axis x)\\
\>\>  (fcn (if (< x axis) x (- (* 2 axis) x)))))
\end{tabbing}\end{tt}


	We can now define our sine expression in terms of these two higher
order expressions:
\begin{tt}\begin{tabbing}
(de\=fine sine\\
\>	(la\=mbda (x)\\
\>\>		(period-maker sine-full (minus (/ pi 2))  (* 2 pi) x)))\\[.10in]
\kill
(de\=fine sine-full\\
\>	(la\=mbda (x)\\
\>\>		(reflection-maker sine-half (/ pi 2) x)))
\end{tabbing}\end{tt}
The function {\tt sine-half}, which computes the value of $\sin(x)$ over
of the interval [$-\pi/2, \pi/2$], is reflected about the axis $\pi/2$
in {\tt sine-full}
so that it can compute $\sin(x)$ for the entire interval [$-\pi/2,
3\pi/2$].

	Now that we have reduced our interval to a reasonable size, we can
turn our attention back to approximating the Taylor series for $\sin(x)$.
The $i$th term in this series can be represented in a direct way 
by {\tt sine-term}:
\begin{tt}\begin{tabbing}
(de\=fine sine-coefficient\\
\>	(la\=mbda (i)\\
\>\>	(/ \=(expt -1.0 i) \\
\>\>\>			(factorial (1+ (* 2 i))))))\\[.10in]
\kill
(de\=fine sine-term \\
\>	(la\=mbda (i x)\\
\>\>		(* \=(sine-coefficient i)\\
\>\>\>		(expt x (1+ (* 2 i)))))
\end{tabbing}\end{tt}
The series can now be approximated by simply adding together 
however many terms we decide are sufficient for our purposes.
Such a series can be generated by the expression {\tt truncated-series}:
\begin{tt}\begin{tabbing}
(de\=fine truncated-series\\
\>	(lam\=bda (term length x)\\
\>\>		(if \=(< length 0)\\
\>\>\>			0.0\\
\>\>\>			(+ \=(term length x)\\
\>\>\>\>				(truncated-series term (1- length) x)))))
\end{tabbing}\end{tt}
Our approximation of $\sin(x)$ using, say, the first 5 terms of the
Taylor series could be completed by defining {\tt sine-half} to be:
\begin{tt}\begin{tabbing}
(de\=fine sine-half\\
\>	(la\=mbda (x)\\
\>\>		(truncated-series sine-term 5 x)))
\end{tabbing}\end{tt}
The expression {\tt (sine pi)} will now reduce to zero with six-place accuracy
in 406 reductions.  

	It is likely that
{\tt sine} will be heavily used in numerical programs, so it is worth
some extra effort to try and cut down on the number of reductions it
takes.  We will now do this in three ways:  
\begin{enumerate}
	\item Change the way the series is represented.
	\item Use only the minimum number of series terms needed to achieve 
the desired accuracy.
	\item Use {\tt rdefine} instead of {\tt define} to `reduce away' much
of the constructive description.
\end{enumerate}


\section{Horner's Rule}

	In the previous section, a simple, straightforward power series was
used to approximate $\sin(x)$.  We can cut down on the number of
arithmetic operations required to compute the series by using Horner's
rule, which rearranges polynomials of the form
\[ P(x) = a_0 + a_1x + a_2x^2 + \cdots + a_nx^n \]
which takes \(n(n+1)/2\) multiplications and $n$ additions to compute,
into the form
\[ P(x) = a_{0} + x(a_1 + x(a_{2} + x(\ldots(a_{n-1} + a_nx)\ldots)))\]
which takes only $n$ multiplications and $n$ additions.
Horner's rule is known to be an optimal way to rearrange a polynomial 
for rapid evaluation without having to do substantial calculations
during the rearrangement.  Horner's rule can be codified as a higher order
expression:
\begin{tt}\begin{tabbing}
(de\=fine horners-rule\\
\>	(la\=mbda (coefficient length variable)\\
\>\>	(let\=rec ((\=term  \\
\>\>\>\>			 	(lam\=bda (n)\\
\>\>\>\>\>				(if \=(= n length)\\
\>\>\>\>\>\>				(coefficient n)\\
\>\>\>\>\>\>				(+ \=(coefficient n)\\
\>\>\>\>\>\>\>					(* variable (term (1+ n))))))))\\
\>\>\>		(term 0))))
\end{tabbing}\end{tt}
{\tt Coefficient} should be a function that produces the
$i$th coefficient for the $i$th term in the series; {\tt length} is the
number of terms in the polynomial.

	The power series we are using to represent $\sin(x)$ is of
the form
\[ a_0x + a_1x^3 + a_2x^5 + \cdots + a_nx^{2n+1} \]
which cannot be directly produced using {\tt horners-rule}.  However, if we
rewrite the series as
\[ x(a_0 + a_1x^2 + a_2x^4 + \cdots + a_nx^{2n}) \]
we can use {\tt horners-rule} to construct the proper series by letting 
{\tt variable} be $x^2$ and multiplying the resulting series by $x$:
\begin{tt}\begin{tabbing}
(de\=fine sine-series\\
\>	(la\=mbda (x x2)\\
\>\>		(* x (horners-rule sine-coefficient 5 x2))))\\[.10in]
\kill
(de\=fine sine-half\\
\> (la\=mbda (x)\\
\>\>	(le\=t ((x2 (* x x))\\
\>\>\>	(sine-series x x2))))
\end{tabbing}\end{tt}
With these changes, {\tt (sine pi)} now reduces to zero in 359
reductions, versus the 406 reductions it previously took.

\section{Minimizing the Number of Terms}

	There is another part of this program to which constructive
programming can be applied, and that is in calculating the number of
terms which need to be evaluated in order to achieve the value of
$\sin(x)$ within a specified accuracy.  In the discussions above, the
number of terms was arbitrarily chosen to be 5.  There is, however, a
way to properly compute this number.

	The absolute error in the truncated sum of an alternating sign,
absolutely convergent series (like the one we are using) is less than
the magnitude of the first term neglected.  Another factor to consider
is that a series should not be truncated until its terms are definitely
getting smaller. For the Taylor series we are using,
\[\sum_{i=0}^{\infty} (-1)^{i} \frac{1}{(2i+1)!} x^{2i+1} \]
the individual terms are monotonically decreasing when the 
denominator starts growing faster
than $x^{2i+1}$, which happens when $2i+1>x$.  Therefore, we must have
at least $\lceil(x-1)/2 \rceil$ terms in our series.

	Putting all of this together yields the expression 
{\tt sine-number-of-terms}, which calculates the number of terms
necessary to achieve accuracy {\tt eps} at the point {\tt value}:

\pagebreak
\begin{tt}\begin{tabbing}
(de\=fine sine-number-of-terms\\
\>   (la\=mbda (eps value)\\
\>\>    (let\=rec ((\=loop\\
\>\>\>\>             (la\=mbda (i)\\
\>\>\>\>\>              (if \=(< (abs ((sine-term i) value)) (abs eps))\\
\>\>\>\>\>\>                 (1- i)\\
\>\>\>\>\>\>                 (loop (1+ i))))))\\
\>\>\>         (loop (ceiling (/ (- value 1) 2))))))
\end{tabbing}\end{tt}
We can now redefine {\tt sine-series} so that it uses {\tt
sine-number-of-terms} instead of hard coding in the number 5.
An additional benefit is that by using different values for {\tt eps}, 
other versions of {\tt sine} can be automatically 
generated which have differing accuracies.


\section{`Reducing Away' the Descriptions}

	Using Horner's Rule and the minimal number of series terms has lowered the
number of reductions (i.e., computational work) required to calculate
$\sin(x)$, but not by a significant amount.  Each time {\tt sine} is
used, it goes through the process of constructing a program to calculate
$\sin(x)$.  To really cut down on the number of reductions, this
repetitive construction cost must be eliminated.

	This is where the strong head normalization semantics of THOR come
into play.  Using the {\tt rdefine} special form, it is possible to
reduce expressions before they are associated with symbols, thus
allowing much of the overhead of constructive description to be `reduced
away' at load time.  This overhead is only paid once, when the
expressions are defined, and not each time {\tt sine} is reduced.

	By selectively choosing the correct {\tt define}'s to change into
{\tt rdefine}'s, the number of reductions required for {\tt sine} can be
cut dramatically.  By {\tt rdefine}-ing {\tt sine}, {\tt sine-full}, and
{\tt sine-series} 
	\footnote{To fully utilize strong head normalization in {\tt
	sine-series}, it is necessary to change it's definition slightly
	so that the {\tt value} parameter of {\tt sine-number-of-terms} is
	not a variable, but a real number.  A good choice for {\tt value}
	is $\pi/2$, which is the axis of reflection for our chosen interval.
	This has only a slight affect on the overall accuracy of {\tt sine}.}
the number of reductions required for {\tt (sine pi)}
is cut to 33.  That's less than one-tenth of the reductions required
previously!

	What exactly is going on here?  When the {\tt redefine}'s are
processed by HORSE, the constructive descriptions build expressions for
computing $\sin(x)$, and these newly constructed expressions are
associated with the {\tt rdefine}'d symbol.  For example, after
rdefinition, the expression associated with {\tt sine-series} is:

\pagebreak
\begin{tt}\begin{tabbing}
(LA\=MBDA (X X2) \\
\>	(* X (+ \=1.000000\\
\>\>	     (* X2 (+ \=-0.166667 \\
\>\>\>			      (* X2 (+ \=0.008333 \\
\>\>\>\>							(* X2 (+ \=-0.000198 \\
\>\>\>\>\>									(* X2 0.000003))))))))))
\end{tabbing}\end{tt}
The constructive definitions of the Taylor series terms, Horner's Rule,
and the minimum number of terms to use have all disappeared, leaving
behind an arithmetic expression.  We can have confidence that this
arithmetic expression is correct, because it was produced from a
mathematically correct description by equivalence preserving transformations.

\section{Discussion}

	A couple of comments need to be made about using {\tt rdefine}.  The
first is that the order of definitions and rdefinitions is important. 
This is nothing new; most programming languages require programs to be
ordered in certain ways for purposes of scoping and such.  
The order in which definitions were presented in this chapter was chosen
deliberately and must be maintained to reap the full benefits of {\tt
rdefine}-ing.   This leads right into the second point, which is
selecting the proper symbols to {\tt rdefine}.

	Selecting the proper symbols to {\tt rdefine} and the
proper ordering for all definitions is partly trial-and-error, partly
art.  For example, if we had
only {\tt rdefine}'d {\tt sine} after everything else had been defined,
{\tt (sine pi)} would take 216 reductions.  An improvement over the
initial situation, but clearly lacking.  A look at the resulting
expression for {\tt sine} reveals that it is huge, littered with
duplicate subexpressions.  It is recomputing all the duplicate subexpressions
that costs so much.  The choice of rdefined symbols presented
above eliminates many of the duplicate expressions, allowing the value
of many arguments to be shared thereby reducing the computation required.

	THOR is a new language with many novel features.  Most of my
energies have been devoted to it's design and implementation;  I have
not had a lot of time to actually write programs in THOR and gain
experience in how to best utilize its new abilities.
Perhaps for Version 2.0 I will have enough experience using THOR to
provide more and better examples and suggestions.


